(1) 波带片平面半径 \(r\) 的点到轴上设计焦点的距离为 \(\sqrt{f^2+r^2}\)，中心点距离为 \(f\)，故相对中心的光程差
\[
\Delta(r)=\sqrt{f^2+r^2}-f.
\]
相邻 Fresnel 半波带边界要求 \(\Delta(r_n)=n\lambda/2\)，\(n=1,2,\ldots\)。解出
\[
r_n^2=n\lambda f+\left(\frac{n\lambda}{2}\right)^2.
\]
在 \(n\lambda\ll f\) 的傍轴范围，第二项相对很小，
\[
r_n\simeq\sqrt{n\lambda f}.
\]
等式中 \(r_n^2\) 与 \(\lambda f\) 都具有长度平方量纲。

(2) 由 Huygens–Fresnel 原理，轴上观察点的复场是各环区域积分的相干和。连续相邻半波带的相位近似相差 \(\pi\)，因此原本相邻环对焦点场有相互抵消的趋势。若把其中一套交替半波带遮挡，剩余透明半波带在设计焦点的主要相位接近同向，从而增强该焦点相对于无结构孔径的相干聚焦效果。透射总功率确实变少，但轴上焦点的局部强度可以增加；两种量不可混淆。

(3) 当波带片刻制的半径 \(r_n\) 不变，在近轴设计关系 \(r_n^2\simeq n\lambda f\) 中，新的工作波长 \(\lambda'\) 对应主焦距
\[
f_1(\lambda')\simeq\frac{r_n^2}{n\lambda'}\propto\frac1{\lambda'}.
\]
因此波长越长，主焦点越靠近波带片，表现出与普通正色散折射透镜相反的色差趋势。按半波带二值周期的傅里叶谐波，理想振幅型波带片除零级透射背景外可形成奇数级聚焦，约在
\[
f_j=f_1/j,\quad j=1,3,5,\ldots
\]
处，其中高阶焦点效率逐渐下降；高阶焦点也须在标量近轴条件下解释。

(4) 把光阑透过率看作占空比 \(1/2\)、周期为两个半波带的二值振幅光栅，第一非零傅里叶谐波的复幅系数模为 \(1/\pi\)，故相对于全部照明入射功率的理想一阶衍射聚焦效率为
\[
\eta_1^{\rm amplitude}=\frac1{\pi^2}\simeq10.1\%.
\]
若交替半波带不遮挡，而分别赋予相差 \(\pi\) 的二值相位延迟，则第一谐波幅度为 \(2/\pi\)，相应理想一阶效率为
\[
\eta_1^{\rm phase}=\frac4{\pi^2}\simeq40.5\%.
\]
这些是理想二值结构在标量衍射模型中的效率，有限环数、材料吸收、制作误差及高数值孔径效应会降低实际焦点品质。
